\section{The original infinitesimal transitions and their obstructions}

\begin{lemma}[Transitions of the genuine infinitesimal exact sequences]
\label{lem:fc-infinitesimal-sequence-transitions}
\leanok
\lean{EnvelopeProperty.FibreInfinitesimalBoundaryTransitions.quotientTransition, EnvelopeProperty.FibreInfinitesimalBoundaryTransitions.quotientTransition_app, EnvelopeProperty.FibreInfinitesimalBoundaryTransitions.sectionMap_transition, EnvelopeProperty.FibreInfinitesimalBoundaryTransitions.quotientTransition_appTop, EnvelopeProperty.FibreInfinitesimalBoundaryTransitions.quotientTransition_baseSectionsMap, EnvelopeProperty.FibreInfinitesimalBoundaryTransitions.sequenceTransition}
\uses{lem:fc-rees-cech-degree-transitions}
For $n\leq m$, the original map $X_n\to X_m$ induces the actual map
of pushforward structure modules on $X$.  On every original open it is
the entire pullback map on the corresponding original inverse-image
open.  Together with the original inclusion
$I^m\mathcal O_X\to I^n\mathcal O_X$ and the identity of
$\mathcal O_X$, it gives a morphism between the two original exact
sequences of infinitesimal structure modules.
\end{lemma}
\begin{proof}
Push the canonical structure-module map for $X_n\to X_m$ forward to
$X$, and use the actual composition equality of the original
projections.  The pushforward composition and equality isomorphisms
give the stated module map; on an open, their definition is restriction
along that composition equality, so its value is exactly the original
scheme pullback on the inverse-image open.  The projection square
commutes because the original transition composed with the original
projection is the original lower-order projection.  The ideal-power
square commutes since each original inclusion retains every underlying
section.  These two squares are the required morphism of exact
sequences.
\end{proof}

\begin{lemma}[Naturality of the original global-section obstruction]
\label{lem:fc-infinitesimal-boundary-transitions}
\leanok
\lean{EnvelopeProperty.FibreInfinitesimalBoundaryTransitions.complexSequenceTransition, EnvelopeProperty.FibreInfinitesimalBoundaryTransitions.connecting_naturality, EnvelopeProperty.FibreInfinitesimalBoundaryTransitions.sectionBoundary_naturality, EnvelopeProperty.FibreInfinitesimalBoundaryTransitions.sectionBoundary_transition_apply}
\uses{lem:fc-infinitesimal-sequence-transitions}
On an ordered affine cover of $X$, the connecting maps to the original
ideal-power cohomology commute with every original infinitesimal
transition.  In degree zero this is naturality of the obstruction of
each entire original infinitesimal global section.
\end{lemma}
\begin{proof}
Apply the original ordered Čech functor to the morphism of exact
sequences.  Its rows are short exact on every affine intersection.
Naturality of the connecting morphism of a short exact sequence of
complexes therefore gives the claimed cohomology square.  In degree
zero, the original global-section to Čech-cohomology isomorphism is
natural for every module morphism.  Composing that naturality square
with the connecting square gives the obstruction square.  The
pushforward transition on entire global sections is the literal
pullback along $X_n\to X_m$, by its whole-open formula.
\end{proof}


\begin{lemma}[The obstruction of a compatible original infinitesimal family]
\label{lem:fc-infinitesimal-compatible-section-boundary}
\leanok
\lean{EnvelopeProperty.FibreInfinitesimalBoundaryTransitions.compatible_sectionBoundary}
\uses{lem:fc-infinitesimal-boundary-transitions,lem:fc-rees-cech-degree-transitions}
Let $U$ be a finite ordered affine cover of the separated $R$-scheme
$X$.  Suppose $(s_n)$ is a compatible family of entire original
sections on the positive-order models $X_{n+1}$.  Identify the
cohomology of $I^{n+1}\mathcal O_X$ with the original degree-$(n+1)$
Rees--\v Cech cohomology using the coefficient-preserving comparison.
Then the degree-one obstructions of $s_n$ form a compatible family in
those same original cohomology modules, for all original
higher-to-lower power maps.
\end{lemma}
\begin{proof}
For $n\leq m$, naturality of the preceding obstruction identifies the
image of the obstruction of $s_m$ under
$H^1(I^{m+1}\mathcal O_X)\to H^1(I^{n+1}\mathcal O_X)$ with the
obstruction of the pullback of $s_m$ to $X_{n+1}$.  The pullback is
$s_n$ by compatibility of the given family.  The comparison with the
original Rees degree complexes intertwines the entire ideal-power
inclusion with its actual higher-to-lower degree map.  Applying that
commuting square gives precisely the asserted equality between the
two compared obstruction classes.  This argument applies to every
$n\leq m$ and uses no vanishing or section-lifting premise.
\end{proof}
